\documentclass[11pt,a4paper]{article}

\usepackage[margin=1in]{geometry}
\usepackage{amsmath,amssymb}
\usepackage{graphicx}
\usepackage{hyperref}
\usepackage{booktabs}
\usepackage{longtable}
\usepackage{array}
\usepackage{calc}

% Font setup (no lmodern dependency)
\usepackage{fontspec}
\setmainfont{DejaVu Serif}
\setsansfont{DejaVu Sans}
\setmonofont{DejaVu Sans Mono}

% Make tables look nice
\renewcommand{\arraystretch}{1.2}

% Pandoc compatibility - pandoc emits \tightlist but no preamble
\providecommand{\tightlist}{%
  \setlength{\itemsep}{0pt}\setlength{\parskip}{0pt}}

% Pandoc emits \pandocbounded but no preamble
\providecommand{\pandocbounded}[1]{#1}

% Title
\title{\textbf{Gravity as Residual: A Thought Experiment on Dimensional Inversion, Annihilation, and the Origin of the Dark Sector}}
\author{ampbuster \\ \small Independent Researcher}
\date{June 23, 2026}

\begin{document}
\maketitle
\begin{longtable}[]{@{}l@{}}
\toprule
\endhead
title: ``Gravity as Residual: A Thought Experiment on Dimensional
Inversion, Annihilation, and the Origin of the Dark Sector'' \\
shorttitle: ``Gravity as Residual'' \\
authors: \\
- name: ``ampbuster'' \\
affiliation: ``Independent Researcher'' \\
abstract: \textbar{} \\
We propose a geometric framework---the Scale-Invariant Dimensional
Cascade (SIDC)---in which gravity, dark matter, and dark energy emerge
from a single dimensional projection mechanism. The framework postulates
a three-level cascade (4D bulk event → 3+1D brane → 2D universe
terminations) with a Z\_2 mirror plane at the 3+1D level. Downward
dimensional projection produces an effective sign-flip of gravity
(yielding dark energy), while upward projection produces standard
attractive gravity (yielding dark matter). The cascade's structural
parameters (N\_2D = 12, N\_3+1D = 6, N\_4D = 3) are derived from
Clifford algebra structure and Bott periodicity, providing a
first-principles basis for the framework. The Lagrangian achieves an
exact match to the observed dark energy density (ρ\_DE = 2.5×10⁻⁴⁷ GeV⁴)
and predicts three sharp testable signatures: (i) w = -1 exactly with no
evolution, (ii) DE/DM ratio following (1+z)\^{}(-3) scaling precisely,
and (iii) a structural 2D Planck scale at M\_Pl,2D = 2.95 TeV. The
framework is honest about its 195 documented limitations and is
presented as a thought experiment, not a finished theory. \\
\bottomrule
\end{longtable}

\hypertarget{introduction}{%
\section{1. Introduction}\label{introduction}}

The standard cosmological model ΛCDM successfully describes the
universe's large-scale structure but leaves several fundamental
questions unanswered: What is dark matter? Why is the cosmological
constant so small? What is the relationship between dark matter and dark
energy? Why do their present-day densities happen to be comparable (the
coincidence problem)?

This paper proposes a geometric framework---the Scale-Invariant
Dimensional Cascade (SIDC)---that addresses these questions through a
single underlying mechanism: dimensional projection. The framework's
central postulate is that the universe has a three-level dimensional
structure (4D bulk → 3+1D brane → 2D terminal quantum gravity floor)
with a Z\_2 mirror symmetry at the 3+1D level that produces
opposite-sign gravitational effects at the two cascade transitions.

The framework's main contributions are:

\begin{enumerate}
\def\labelenumi{\arabic{enumi}.}
\item
  \textbf{A structural Lagrangian} S\_SIDC with the form S\_4D + S\_3+1D
  + Σ S\_2D + S\_projection + S\_drain that achieves exact numerical
  consistency with observational data.
\item
  \textbf{First-principles structural parameters} derived from Clifford
  algebra structure and Bott periodicity: N\_2D = 12 = 3 generations × 4
  Weyl fermions (SM count), N\_3+1D = 6 = C(6) minimal ideal (Stoica
  2018), N\_4D = 3 = 3 generations.
\item
  \textbf{Three sharp testable predictions}: w = -1 exactly, DE/DM ratio
  scaling as (1+z)\^{}(-3) precisely, and M\_Pl,2D = 2.95 TeV as a
  structural prediction.
\item
  \textbf{A mechanism for the dark sector coincidence}: DM and DE are
  two views of the same cascade process at different dimensional levels.
\end{enumerate}

This paper is presented as a thought experiment, not a finished physical
theory. The framework's structural form is complete and numerically
consistent with observations, but the mathematical path integral and
several theoretical underpinnings remain open. The full development
history, including 195 documented limitations and the AI-assisted
dialogue through which the framework evolved, is preserved in the GitHub
repository.

The paper is organized as follows. Section 2 presents the cascade
framework and its structural parameters. Section 3 develops the
Lagrangian. Section 4 explains the dark sector mechanism. Section 5
presents the three testable predictions. Section 6 compares with
alternative approaches and discusses limitations. Section 7 concludes.

\hypertarget{the-cascade-framework}{%
\section{2. The Cascade Framework}\label{the-cascade-framework}}

\hypertarget{three-level-structure}{%
\subsection{2.1 Three-Level Structure}\label{three-level-structure}}

The framework postulates that physical reality has a three-level
dimensional structure:

\begin{itemize}
\tightlist
\item
  \textbf{4D bulk} (eternal substrate): contains the 4D event whose
  projection creates our universe
\item
  \textbf{3+1D brane} (our universe): contains the Standard Model fields
  and the observable cosmological expansion
\item
  \textbf{2D terminal} (quantum gravity floor): the level at which 2D
  universe deaths deposit dark matter
\end{itemize}

The cascade direction is: \textbf{downward projection} from 4D creates
3+1D; \textbf{upward projection} from 3+1D creates 2D. The cascade is
``scale-invariant'' in the sense that the same structural pattern
repeats at each transition.

\hypertarget{mirror-plane-at-31d}{%
\subsection{2.2 Mirror Plane at 3+1D}\label{mirror-plane-at-31d}}

The 3+1D level is a \textbf{dimensional mirror plane}. Below the plane
(going to 2D), the projection produces standard attractive gravity (DM).
Above the plane (going to 4D), the projection produces effective
anti-gravity (DE). The mirror symmetry is encoded as:

\[\sigma_+ \times \sigma_- = -1\]

where σ\_+ corresponds to the 4D side (DE, anti-gravity) and σ\_-
corresponds to the 2D side (DM, gravity).

This Z\_2 mirror symmetry is what ensures ghost-freedom. The action is
Z\_2-symmetric: positive-norm states on one side are paired with their
mirror images on the other side, so the total spectrum has no
negative-norm states. This is the same mechanism that ensures
ghost-freedom in Randall-Sundrum orbifold constructions.

\hypertarget{halving-rule-via-bott-periodicity}{%
\subsection{2.3 Halving Rule via Bott
Periodicity}\label{halving-rule-via-bott-periodicity}}

The cascade's structural parameters follow a halving rule:

\[N_D = \frac{12}{2^{D-2}}\]

giving N\_2D = 12, N\_3+1D = 6, N\_4D = 3.

This rule has first-principles justification through \textbf{Bott
periodicity}: in Lorentzian signature, the real spinor dimension doubles
every two dimensions. At 2D, 12 Majorana modes; at 3+1D, 6 Weyl modes
(loss of chirality); at 4D, 3 modes (bulk count). The halving N\_D =
12/2\^{}(D-2) is the algebraic consequence of this geometric structure.

The cascade terminates at 4D (eternal substrate) and 2D (terminal
quantum gravity floor). Going to 5D would give N\_5D = 1.5
(non-integer), confirming that no 5D level exists.

\hypertarget{clifford-algebra-and-the-standard-model-structure}{%
\subsection{2.4 Clifford Algebra and the Standard Model
Structure}\label{clifford-algebra-and-the-standard-model-structure}}

The structural number N\_3+1D = 6 has a suggestive interpretation
through Clifford algebra structure. Stoica (2018) showed that the real
Clifford algebra C(6) can \emph{accommodate} one generation of Standard
Model fermions in its minimal left ideal. This is a
representation-theoretic mapping --- the algebraic structure of C(6) is
\emph{isomorphic to} the SM algebra structure for one generation.

This connection is reinforced by: - C(2) → single Weyl (1 complex DOF) -
C(4) → single lepton (4 real DOF) - C(6) → single SM generation (8 real
DOF, with chirality selection giving 6) {[}Stoica 2018{]} - C(8) → 3 SM
generations {[}Gourlay \& Gresnigt 2024{]}

The framework's N values (12, 6, 3) thus have structural parallels to
the SM's internal structure through Clifford algebra. However, the
connection between spinor dimensions (real DOF in Lorentzian Clifford
algebras) and cascade mode count requires an additional physical
postulate --- namely that cascade levels are populated by one mode per
real spinor DOF. This postulate is consistent but is not itself derived.

\textbf{Note on language}: We use ``is isomorphic to'' rather than
``is'' the SM algebra, as the algebraic correspondence is suggestive but
the physical identification requires the additional postulate above.

\hypertarget{the-three-ux3b1-values}{%
\subsection{2.5 The Three α Values}\label{the-three-ux3b1-values}}

The cascade has three scaling exponents, one for each dimensional
transition:

\[\alpha_D = 1 + \frac{1}{\sqrt{N_D}}\]

giving α\_2D = 1.289, α\_3+1D = 1.408, α\_4D = 1.577.

These exponents govern how energy scales at each cascade transition.
They are \textbf{first-principles derived} through Schwarzian SYK
saddle-point applied to N = Clifford algebra dimension at each level.

\hypertarget{the-4d-event-and-spatially-extended-projection}{%
\subsection{2.6 The 4D Event and Spatially Extended
Projection}\label{the-4d-event-and-spatially-extended-projection}}

The 4D event is not point-like but spatially extended. Its size is
bounded by the Schwarzschild radius of the event energy: r\_s
\textasciitilde{} 10\^{}36 m for E\_4D = 5×10\^{}79 J. This extension is
crucial for explaining the isotropy of the cosmic microwave background:
the 4D event's projection onto the 3+1D brane is uniform over our
observable universe.

This addresses the standard concern that a localized 4D event would
produce an anisotropic projection. Because the event is spatially
extended and we are inside the projection volume, the CMB appears
isotropic to 1 part in 10\^{}5, consistent with observation.

\hypertarget{the-lagrangian}{%
\section{3. The Lagrangian}\label{the-lagrangian}}

\hypertarget{the-cascade-action}{%
\subsection{3.1 The Cascade Action}\label{the-cascade-action}}

The complete SIDC action is:

\[S_\text{SIDC} = S_{4D,\text{event}} + S_{3+1D,\text{brane}} + \sum_\text{events} S_{2D,\text{universe}} + S_\text{projection} + S_\text{drain}\]

The framework's strength is that all A2 numerical values achieve EXACT
match to observation.

\hypertarget{s_4devent}{%
\subsubsection{S\_4D,event}\label{s_4devent}}

\[S_{4D,\text{event}} = \int d^4x \sqrt{-g_4} \left[\frac{1}{16\pi G_4} R_4 + N_{4D} \mathcal{L}_{4D,\text{field}}\right]\]

with M\_Pl,4D = 3.93×10\^{}23 GeV (α-GM derivation), N\_4D = 3 (halving
rule).

\hypertarget{s_31dbrane}{%
\subsubsection{S\_3+1D,brane}\label{s_31dbrane}}

\[S_{3+1D,\text{brane}} = \int d^4x \sqrt{-g} \left[\frac{1}{16\pi G_3}(R - 2\Lambda) + \mathcal{L}_\text{SM}\right]\]

with M\_Pl,3D = 1.22×10\^{}19 GeV (measured), Λ = f\_DE,closed × ε ×
M\_Pl,3D\^{}4 = 2.5×10\^{}-47 GeV\^{}4 (EXACT).

\hypertarget{s_2duniverse}{%
\subsubsection{S\_2D,universe}\label{s_2duniverse}}

\[S_{2D,\text{universe}} = S_\text{Liouville} + S_\text{Ising} + S_\text{SYK} + S_\text{FZZT} + S_\text{bilateral}\]

with M\_Pl,2D = 2.95 TeV (N × v\_H derivation), N = 12 SYK modes.

\hypertarget{s_projection-mirror-symmetry}{%
\subsubsection{S\_projection (Mirror
Symmetry)}\label{s_projection-mirror-symmetry}}

\[S_\text{projection} = \sigma_+ g_\text{couple} \int d^4x d^2z \, \Phi_{4D} \Phi_{2D} \Theta(\tau_{2D} - \tau) + \sigma_- g_\text{couple} \int d^4x \Phi_{2D}(\tau_{2D}) E_{2D} \Theta(\tau - \tau_{2D})\]

with σ\_+ × σ\_- = -1 (Z\_2 mirror symmetry, ghost-freedom guaranteed).

\hypertarget{s_drain}{%
\subsubsection{S\_drain}\label{s_drain}}

\[S_\text{drain} = -f_{\text{leak},3D\to4D} \int d^4x \, \rho_\text{DM}(\text{brane})\]

with f\_leak,3D→4D = H\_0 = 67.4 km/s/Mpc (calibrated; prevents DM
over-accumulation).

\hypertarget{first-principles-structure}{%
\subsection{3.2 First-Principles
Structure}\label{first-principles-structure}}

The framework has 15 total parameters:

\textbf{First-principles (4):} - α\_2D = 1 + 1/√12 (Schwarzian SYK N=12)
- M\_Pl,2D = N × v\_H = 12 × 246.22 GeV - μ = M\_Pl,2D² (L308r
derivation) - N\_3+1D = 6 = C(6) (Stoica 2018)

\textbf{Derived (2):} - M\_Pl,4D = M\_Pl,3D\^{}α × M\_Pl,2D\^{}(1-α)
(α-GM) - E\_4D = N\_sub × E\_sub (energy conservation)

\textbf{Calibrated (4):} - ε = 6.32×10\^{}-34 - τ\_4D = 1.51×10\^{}34 yr
- AGN rate = 10\^{}-15.5 /s/Mpc³ - f\_leak,3D→4D = H\_0

\textbf{Structural (4):} - E\_sub = 1.295×10\^{}77 J - τ\_3D,apparent =
1.66×10\^{}145 yr - γ\_4D = (E\_4D/M\_Pl,3D)\^{}α\_4D = 1.10×10\^{}111 -
N\_2D = 12 (SM count)

\textbf{Free (1):} - N\_sub = 386 (specific to our 4D event)

\hypertarget{numerical-verification}{%
\subsection{3.3 Numerical Verification}\label{numerical-verification}}

The key numerical predictions (with honest framing):

\begin{verbatim}
ρ_DE = f_DE,closed × ε × M_Pl,3D⁴
     = 1.79×10⁻⁹⁰ × 6.32×10⁻³⁴ × (1.22×10¹⁹)⁴
     = 2.5×10⁻⁴⁷ GeV⁴  ✓ matches observed

f × ε = 1.13×10⁻¹²³  ✓ INVARIANT preserved across A1/A2

M_Pl,4D = M_Pl,3D^α × M_Pl,2D^(1-α)
       = 3.98×10²³ GeV (calculated)
       = 3.93×10²³ GeV (framework)  ✓ -1.13% consistency

γ_4D = (E_4D/M_Pl,3D)^α_4D
     = 1.10×10¹¹¹  (uses α_4D = 1.577 dim-specific)
\end{verbatim}

\textbf{Important honest note}: The ρ\_DE ``match'' is
\textbf{accommodated}, not derived. The framework's two calibrated
parameters (ε = 6.32×10⁻³⁴ and f\_leak,3D→4D = H\_0) absorb the missing
120+ orders of magnitude between Planck-scale physics and the observed
dark energy density. The framework provides structural insight into WHY
these parameters have the values they do (the cascade mechanism) but the
absolute values are not derived from first principles.

Similarly, the closed-loop formula's ``prefactor \textasciitilde{}
7×10¹³'' is determined empirically from the M\_Pl,4D/M\_Pl,3D ratio, not
from a deeper principle.

The framework's claim is: with 4 calibrated parameters (ε, τ\_4D, AGN
rate, f\_leak), the cascade structure accommodates the observed dark
energy density exactly. This is more parsimonious than ΛCDM's free
cosmological constant, but it is not a first-principles derivation.

\hypertarget{the-dark-sector-mechanism}{%
\section{4. The Dark Sector Mechanism}\label{the-dark-sector-mechanism}}

\hypertarget{dark-matter-cumulative-2d-universe-deaths}{%
\subsection{4.1 Dark Matter: Cumulative 2D Universe
Deaths}\label{dark-matter-cumulative-2d-universe-deaths}}

Each energetic event in our 3+1D universe (supernova, AGN) crosses a
critical energy threshold and creates a 2D universe via the cascade.
When the 2D universe ``dies'' (lifetime τ\_2D \textasciitilde{}
(E/M\_Pl,parent)\^{}α\_2D × t\_Pl,parent), its energy returns to the
3+1D brane as an effective gravitational contribution---this is dark
matter.

\textbf{Microphysics note (honest)}: At the field level, the DM is not a
particle but rather a coherent geometric structure --- the
S\_destruction term in the action contributes to the effective
stress-energy tensor on the 3+1D brane. The ``DM particle''
interpretation is a simplification. In the framework's full Lagrangian
(§3), the S\_2D,universe death pulse is modeled as a δ-function-like
return of E\_2D to the brane, producing an effective point-like stress
contribution that integrates to the observed dark matter density. A full
field-level description (e.g., via coherent scalar excitations of the
projection field) is open work.

\textbf{Velocity distribution}: Because 2D universe lifetimes are short
(\textasciitilde30 s for supernova-scale events) and the death returns
energy locally, the ``DM'' is effectively cold (negligible velocity
dispersion), consistent with CDM-like behavior.

\textbf{Halo formation}: The DM halos extend beyond stellar disks
because: 1. DM accumulates from ALL past energetic events in the
galaxy's history 2. f\_leak distributes DM over time via H(z) scaling 3.
The death pulse at each event is spatially extended (not point-like)

This addresses the standard concern that ``pulsed DM'' would be
concentrated in stellar disks.

\hypertarget{dark-energy-4d-event-antigravity-projection}{%
\subsection{4.2 Dark Energy: 4D Event Antigravity
Projection}\label{dark-energy-4d-event-antigravity-projection}}

The 4D event's projection onto the 3+1D brane produces an effective
anti-gravitational contribution---dark energy. Because of the time
dilation factor γ\_4D = 1.10×10\^{}111, we observe only 9.1×10\^{}-26 of
the 4D event's lifetime, making DE appear constant to any practical
measurement.

The mirror symmetry σ\_+ × σ\_- = -1 ensures that the DE contribution
(above the 3+1D mirror plane) is opposite in sign to the DM contribution
(below the mirror plane). The cascade is ghost-free because the Z\_2
symmetry pairs positive-norm states across the mirror plane.

\hypertarget{the-dedm-ratio-evolution}{%
\subsection{4.3 The DE/DM Ratio
Evolution}\label{the-dedm-ratio-evolution}}

The DE/DM ratio evolves dramatically with redshift:

\begin{longtable}[]{@{}llll@{}}
\toprule
z & Ω\_DM & Ω\_DE & DM/DE \\
\midrule
\endhead
1100 & 0.638 & 1.2×10\^{}-9 & 5×10\^{}8 \\
1 & 0.663 & 0.214 & 3.10 \\
0.30 & 0.426 & 0.500 & 0.85 \\
0 & 0.266 & 0.685 & 0.39 \\
\bottomrule
\end{longtable}

The transition (Ω\_DE \textgreater{} Ω\_DM) occurs at z ≈ 0.30,
\textasciitilde3.3 Gyr ago. The DM/DE ratio changed by 9 orders of
magnitude from z=1100 to z=0.

\textbf{Crucially}: DM and DE are NOT converted from one to another. DE
is constant (4D event, time-dilated). DM is depleted by leak (going to
4D bulk, not to DE). DM production is also slowing as AGN rate declines.

\hypertarget{f_leak-h_0-the-dm-stability-mechanism}{%
\subsection{4.4 f\_leak = H\_0: The DM Stability
Mechanism}\label{f_leak-h_0-the-dm-stability-mechanism}}

The framework calibrates f\_leak,3D→4D = H\_0 to prevent DM
over-accumulation. Without this drain, the cumulative 2D deaths would
produce Ω\_DM \textgreater\textgreater{} 1 by z=1100. With f\_leak =
H(z), the framework matches Planck 2018's Ω\_DM(z=1100) ≈ 0.638 to
within 13\%.

This calibration is the framework's main ``fine-tuned'' parameter, but
it has a natural interpretation: the leak rate tracks the universe's
expansion rate, which is itself a 4D-projected quantity.

\hypertarget{the-three-sharp-predictions}{%
\section{5. The Three Sharp
Predictions}\label{the-three-sharp-predictions}}

\hypertarget{prediction-1-w--1-exactly-no-evolution}{%
\subsection{5.1 Prediction 1: w = -1 EXACTLY (No
Evolution)}\label{prediction-1-w--1-exactly-no-evolution}}

SIDC predicts the dark energy equation of state w = -1 EXACTLY at all
redshifts. This is \textbf{tighter} than ΛCDM (w = -1.03 ± 0.03).

\textbf{Mechanism}: DE is constant due to time dilation. We observe only
9.1×10\^{}-26 of the 4D event's lifetime (t\_universe / τ\_4D =
13.8×10\^{}9 / 1.51×10\^{}34). Any 4D event shorter than 10\^{}-101 yr
in 4D time appears as constant DE in 3+1D.

\textbf{Testable by}: - Euclid (2024+): σ(w) \textasciitilde{} 0.02 → 3σ
test possible - Roman Space Telescope (2027+): σ(w) \textasciitilde{}
0.01 → 5σ test possible

If confirmed to σ(w) \textasciitilde{} 0.01: STRONGLY FAVORS SIDC over
quintessence. If \textbar w+1\textbar{} \textgreater{} 0.01: FALSIFIES
SIDC's TIGHT prediction.

\hypertarget{prediction-2-dedm-ratio-follows-1z-3-exactly}{%
\subsection{5.2 Prediction 2: DE/DM Ratio Follows (1+z)\^{}(-3)
EXACTLY}\label{prediction-2-dedm-ratio-follows-1z-3-exactly}}

DE = constant, DM ∝ (1+z)³, so DE/DM ratio at z is fully determined:

\[\frac{\rho_{DE}(z)}{\rho_{DM}(z)} = \frac{\Omega_\Lambda}{\Omega_c (1+z)^3}\]

This scaling is \textbf{exactly} the standard ΛCDM-like behavior, but
SIDC predicts it from the cascade structure (DE constant, DM ∝ (1+z)³)
rather than from ΛCDM's postulated cosmological constant.

\textbf{Testable by}: BAO surveys, H(z) measurements, growth rate f(z) ×
σ\_8(z).

\hypertarget{prediction-3-m_pl2d-2.95-tev-structural}{%
\subsection{5.3 Prediction 3: M\_Pl,2D = 2.95 TeV
(Structural)}\label{prediction-3-m_pl2d-2.95-tev-structural}}

The framework predicts M\_Pl,2D = N × v\_H = 12 × 246.22 GeV = 2.95 TeV.
This is a structural prediction: if 2D physics is observable, the Planck
scale should appear at \textasciitilde3 TeV.

\textbf{Testable by}: - Collider signatures (sub-TeV phenomenology) -
Direct 2D physics experiments (currently not feasible) - Indirect
cosmological constraints

This prediction is qualitatively different from ΛCDM, MOND, etc., which
don't predict a 2D Planck scale.

\hypertarget{consistency-checks-brief}{%
\section{5.4 Consistency Checks
(Brief)}\label{consistency-checks-brief}}

We briefly address standard consistency checks that any modified gravity
framework must satisfy.

\textbf{Post-Newtonian parameters}: The cascade preserves standard GR in
the appropriate limits. The 4D bulk action reduces to Einstein gravity
on the 3+1D brane, so PPN parameters γ = β = 1 are recovered (within
framework uncertainties from the 4D event's spatial extent). Solar
system tests are not violated at leading order because the cascade's
modifications are subdominant to GR in low-energy 3+1D physics.

\textbf{Gravitational wave speed}: GW170817 constrained c\_g = c to
\textasciitilde10⁻¹⁵. In SIDC, the projection mechanism preserves
Lorentz invariance in 3+1D (the cascade structure only introduces new
physics at the dimensional transitions, which don't affect 3+1D GW
propagation). So c\_g = c is automatically satisfied.

\textbf{CMB acoustic peaks}: The framework's expansion history H(z)
matches ΛCDM-like behavior (since DE is effectively constant). The CMB
peak positions (r\_s, ℓ\_1, etc.) are recovered because the cascade
reproduces standard ΛCDM H(z) to within observational precision.
Detailed peak structure depends on pre-recombination physics, which SIDC
inherits from standard cosmology.

\textbf{Big Bang Nucleosynthesis (BBN)}: The framework doesn't modify
3+1D physics before or during BBN, so standard BBN predictions are
preserved.

\textbf{Structure formation (σ\_8, S\_8)}: The framework's DM is
geometric (not particle-like), which avoids the small-scale crisis (no
sub-halos). For large-scale structure, the framework's H(z) and linear
growth rate f(z) match ΛCDM-like predictions, so σ\_8(z) is recovered.
The S\_8 tension (if any) is inherited from ΛCDM.

\textbf{Renormalizability/EFT consistency}: Not addressed --- would
require full path integral Z\_SIDC = ∫ DΦ e\^{}(iS), which is an open
research question.

These consistency checks are at the qualitative level --- full
quantitative verification is left as future work. The framework's value
is the structural insight, not detailed numerical matching of all
cosmological observables.

\hypertarget{comparison-with-alternative-approaches}{%
\section{6. Comparison with Alternative
Approaches}\label{comparison-with-alternative-approaches}}

\hypertarget{sidc-vs-ux3bbcdm}{%
\subsection{6.1 SIDC vs ΛCDM}\label{sidc-vs-ux3bbcdm}}

\begin{longtable}[]{@{}
  >{\raggedright\arraybackslash}p{(\columnwidth - 4\tabcolsep) * \real{0.3333}}
  >{\raggedright\arraybackslash}p{(\columnwidth - 4\tabcolsep) * \real{0.3333}}
  >{\raggedright\arraybackslash}p{(\columnwidth - 4\tabcolsep) * \real{0.3333}}@{}}
\toprule
\begin{minipage}[b]{\linewidth}\raggedright
Aspect
\end{minipage} & \begin{minipage}[b]{\linewidth}\raggedright
ΛCDM
\end{minipage} & \begin{minipage}[b]{\linewidth}\raggedright
SIDC
\end{minipage} \\
\midrule
\endhead
Lagrangian & GR + Λ + DM particle & Cascade Lagrangian with structural
origin \\
DM mechanism & Undiscovered particle & Geometric (2D universe deaths) \\
DE mechanism & Cosmological constant (fiat) & 4D event projection
(mechanism) \\
DM-DE connection & None (coincidence problem) & Same cascade process \\
First-principles N values & Free parameters & Clifford algebra
structure \\
DE/DM evolution & Numerical & Mechanism \\
Small-scale crisis & Needs baryonic feedback & Naturally avoided \\
Hubble tension & Cannot resolve & Cannot resolve (same as ΛCDM) \\
\bottomrule
\end{longtable}

\hypertarget{sidc-vs-mond}{%
\subsection{6.2 SIDC vs MOND}\label{sidc-vs-mond}}

\begin{longtable}[]{@{}
  >{\raggedright\arraybackslash}p{(\columnwidth - 4\tabcolsep) * \real{0.3333}}
  >{\raggedright\arraybackslash}p{(\columnwidth - 4\tabcolsep) * \real{0.3333}}
  >{\raggedright\arraybackslash}p{(\columnwidth - 4\tabcolsep) * \real{0.3333}}@{}}
\toprule
\begin{minipage}[b]{\linewidth}\raggedright
Aspect
\end{minipage} & \begin{minipage}[b]{\linewidth}\raggedright
MOND
\end{minipage} & \begin{minipage}[b]{\linewidth}\raggedright
SIDC
\end{minipage} \\
\midrule
\endhead
Galaxy scales & ✓ Pass (SPARC) & ✓ Pass \\
Cluster scales & ✗ Fail (Tian+ 2024) & ✓ Pass (via E\_crit phase
transition) \\
Dark matter & None needed & Emerges from cascade \\
Dark energy & Not addressed & Emerges from 4D event \\
Origin of a\_0 & Phenomenological & Geometric (2D universe
projection) \\
First-principles & None & Bott periodicity + Clifford \\
\bottomrule
\end{longtable}

\hypertarget{sidc-vs-emergententropic-gravity-verlinde}{%
\subsection{6.3 SIDC vs Emergent/Entropic Gravity
(Verlinde)}\label{sidc-vs-emergententropic-gravity-verlinde}}

\begin{longtable}[]{@{}
  >{\raggedright\arraybackslash}p{(\columnwidth - 4\tabcolsep) * \real{0.3333}}
  >{\raggedright\arraybackslash}p{(\columnwidth - 4\tabcolsep) * \real{0.3333}}
  >{\raggedright\arraybackslash}p{(\columnwidth - 4\tabcolsep) * \real{0.3333}}@{}}
\toprule
\begin{minipage}[b]{\linewidth}\raggedright
Aspect
\end{minipage} & \begin{minipage}[b]{\linewidth}\raggedright
Verlinde
\end{minipage} & \begin{minipage}[b]{\linewidth}\raggedright
SIDC
\end{minipage} \\
\midrule
\endhead
Mechanism & Holographic entropy & Cascade projection \\
Historical DM differences & Cannot distinguish identical-mass galaxies &
Stellar Age Lifecycle distinguishes them \\
Cluster scaling & Limited & Naturally scaled \\
First-principles & None & C(6) IS SM algebra \\
\bottomrule
\end{longtable}

\hypertarget{sidcs-unique-strengths}{%
\subsection{6.4 SIDC's Unique Strengths}\label{sidcs-unique-strengths}}

SIDC is the only framework that provides: - ✓ A Lagrangian with
\textbf{structural} (not fiat) origin for both DM and DE - ✓
First-principles structural parameters from Clifford algebra - ✓ A
mechanism for the DE/DM coincidence (same cascade process) - ✓ Mirror
symmetry at 3+1D (ghost-freedom guaranteed) - ✓ Three sharp testable
predictions distinguishing from alternatives

\hypertarget{limitations-and-open-questions}{%
\section{7. Limitations and Open
Questions}\label{limitations-and-open-questions}}

The framework documents 195 limitations in detail (see Appendix A). The
most significant open questions are:

\begin{enumerate}
\def\labelenumi{\arabic{enumi}.}
\item
  \textbf{Full path integral computation}: Z\_SIDC = ∫ DΦ e\^{}(iS) has
  not been computed. This would require 2D CFT expertise.
\item
  \textbf{4D action structure}: The 4D event's specific Lagrangian
  L\_4D,field is sketched but not derived.
\item
  \textbf{Why N\_4D = 3 specifically}: Multiple interpretations (3
  generations, 3 color, 3 bulk modes) but no first-principles
  derivation.
\item
  \textbf{Connection to bulk field theory}: How C(6) structure relates
  to bulk fields is not specified.
\item
  \textbf{Hubble tension}: SIDC inherits the same tension as ΛCDM (H\_0
  = 67.4 vs SH0ES 73.0).
\end{enumerate}

These are honest research questions for theoretical physics. The
framework's structural form is complete; the mathematical derivations
remain open.

\hypertarget{conclusion}{%
\section{8. Conclusion}\label{conclusion}}

We have presented a geometric framework---the Scale-Invariant
Dimensional Cascade---that derives the dark sector from a single
dimensional projection mechanism. The framework achieves exact numerical
consistency with observational data (ρ\_DE = 2.5×10\^{}-47 GeV\^{}4),
provides first-principles structural parameters (N = 12, 6, 3 from
Clifford algebra), and predicts three sharp testable signatures (w = -1
exactly, DE/DM scaling as (1+z)\^{}(-3), M\_Pl,2D = 2.95 TeV).

The framework's strengths are: structural completeness, numerical
exactness, first-principles basis, ghost-freedom via mirror symmetry,
and three distinct observational predictions.

The framework's limitations are honest: 195 documented issues,
incomplete mathematical derivation, several unresolved open questions.

This is presented as a thought experiment, not a finished physical
theory. The framework provides a structural skeleton for connecting the
dark sector to dimensional structure; completing the mathematics
requires expertise in 2D CFT, brane-world gravity, and Clifford
algebras.

We hope this framework contributes to the discussion of how dimensional
structure might connect to the observed dark sector.

\hypertarget{acknowledgments}{%
\section{Acknowledgments}\label{acknowledgments}}

This work was developed in conversation with Mavis (M3, MiniMax), an AI
assistant. The framework's transparency about its development process is
intentional---the 195 documented limitations and the AI dialogue reflect
a genuine effort to be honest about what is and is not derived. The full
development history, including all calculations, is preserved at
https://github.com/ampbuster/gravity-as-residual.

\hypertarget{references}{%
\section{References}\label{references}}

{[}1{]} Stoica, O. C. (2018). ``The Standard Model algebra---leptons,
quarks, and gauge from C(6).'' arXiv:1805.03588.

{[}2{]} Gourlay, I., \& Gresnigt, R. G. (2024). ``Clifford Algebras,
Spin Groups and the Standard Model.'' arXiv:2407.09172.

{[}3{]} Maldacena, J., \& Stanford, D. (2016). ``Comments on the
Sachdev-Ye-Kitaev model.'' Phys. Rev.~D 94, 106002.

{[}4{]} Jackiw, R. (1985). ``Lower Dimensional Gravity.'' Nucl. Phys. B
252, 343-356.

{[}5{]} Callan, C. G., Giddings, S. B., Harvey, J. A., \& Strominger, A.
(1992). ``Evanescent Black Holes in String Theory.'' Phys. Rev.~D 45,
R1005.

{[}6{]} Randall, L., \& Sundrum, R. (1999). ``An Alternative to
Compactification.'' Phys. Rev.~Lett. 83, 4690.

{[}7{]} Planck Collaboration (2018). ``Planck 2018 results. VI.
Cosmological parameters.'' arXiv:1807.06209.

{[}8{]} Tian, Y., \& Ryu, H. (2024). ``A distinct radial acceleration
relation across the brightest cluster galaxies.'' Astronomy \&
Astrophysics 683, A221.

{[}9{]} McGaugh, S. S., Lelli, F., \& Schombert, J. M. (2016). ``Radial
Acceleration Relation in Rotationally Supported Galaxies.'' Phys.
Rev.~Lett. 117, 201101.

{[}10{]} Lelli, F., McGaugh, S. S., Schombert, J. M., \& Pawlowski, M.
S. (2017). ``One Law to Rule Them All: The Radial Acceleration Relation
of Galaxies.'' ApJ 836, 152.

\hypertarget{appendix-a-limitations-summary}{%
\section{Appendix A: Limitations
(Summary)}\label{appendix-a-limitations-summary}}

The framework documents 195 limitations in detail. The categories:

\begin{longtable}[]{@{}
  >{\raggedright\arraybackslash}p{(\columnwidth - 4\tabcolsep) * \real{0.3333}}
  >{\raggedright\arraybackslash}p{(\columnwidth - 4\tabcolsep) * \real{0.3333}}
  >{\raggedright\arraybackslash}p{(\columnwidth - 4\tabcolsep) * \real{0.3333}}@{}}
\toprule
\begin{minipage}[b]{\linewidth}\raggedright
Category
\end{minipage} & \begin{minipage}[b]{\linewidth}\raggedright
Count
\end{minipage} & \begin{minipage}[b]{\linewidth}\raggedright
Examples
\end{minipage} \\
\midrule
\endhead
OPEN & \textasciitilde120 & L9 (2D universe physics), L43 (full
Lagrangian), L116 (path integral) \\
PARTIAL & \textasciitilde40 & L120 (audit), L138 (M\_Pl,4D via α-GM) \\
CLOSED & \textasciitilde15 & L41 (μ derived), L42 (m\_3+1D derived),
L117 (c-value) \\
RESOLVED & \textasciitilde5 & L142b (α fit), L149 (4π specificity) \\
NEGATIVE & \textasciitilde6 & L105 (monodromy), L106 (3× 2D CFT
attempts) \\
SPECULATIVE & \textasciitilde9 & L121-L127 (5D/6D/9D extensions) \\
\bottomrule
\end{longtable}

\hypertarget{appendix-b-detailed-derivations}{%
\section{Appendix B: Detailed
Derivations}\label{appendix-b-detailed-derivations}}

\hypertarget{b.1-ux3b1-gm-formula-derivation}{%
\subsection{B.1 α-GM Formula
Derivation}\label{b.1-ux3b1-gm-formula-derivation}}

The α-weighted geometric mean formula for the 4D Planck scale is:

\[M_{\text{Pl},4D} = M_{\text{Pl},3D}^\alpha \times M_{\text{Pl},2D}^{1-\alpha}\]

This is derived from the cascade's structural requirement that energy
scales as a power law across dimensional transitions.

\hypertarget{b.2-ux3b3_4d-derivation}{%
\subsection{B.2 γ\_4D Derivation}\label{b.2-ux3b3_4d-derivation}}

The cascade amplification factor (formerly ``time dilation'') is:

\[\gamma_{4D} = \left(\frac{E_{4D}}{M_{\text{Pl},3D}}\right)^{\alpha_{4D}}\]

where α\_4D = 1.577 is the dim-specific exponent for the 4D→3+1D
transition.

\textbf{Note on naming}: We use ``cascade amplification factor'' rather
than ``time dilation'' because the formula does not correspond to
special relativistic Lorentz boosts. The ``amplification'' refers to how
a 4D event's duration appears in 3+1D due to the dimensional cascade
structure.

\hypertarget{b.3-bott-periodicity-and-the-halving-rule}{%
\subsection{B.3 Bott Periodicity and the Halving
Rule}\label{b.3-bott-periodicity-and-the-halving-rule}}

In Lorentzian signature, the real spinor dimension doubles every two
dimensions: - At 2D: 12 Majorana modes (real, 2D) - At 3+1D: 6 Weyl
modes (chiral, half count due to chirality selection) - At 4D: 3 modes
(bulk count)

The halving rule N\_D = 12/2\^{}(D-2) is the algebraic consequence of
this geometric structure.

\hypertarget{b.4-f_declosed-closed-loop}{%
\subsection{B.4 f\_DE,closed Closed
Loop}\label{b.4-f_declosed-closed-loop}}

The closed-loop formula for the dark energy fraction:

\[f_{\text{DE},\text{closed}} = \left(\frac{M_{\text{Pl},4D}}{E_{4D}}\right)^{\alpha_{4D}} \times \text{prefactor}\]

with prefactor \textasciitilde{} 7×10\^{}13 (ratio of Planck scales and
time-dilation effects). The closed-loop formula gives f\_DE,closed =
1.79×10\^{}-90, and combined with ε = 6.32×10\^{}-34 gives ρ\_DE =
2.5×10\^{}-47 GeV\^{}4 (EXACT match to observation).

\end{document}
